\section*{数学核验的范围}
数值核验采用有限词表的确定性求和，对照解析恒等式、自动微分和独立的概率公式。重点包括重心超额损失、$J$ 的梯度与 JS 端点、四种 inner 的梯度和 Hessian、成对估计的协方差、形状端点、TOP-D 算术收缩以及条件与序列的反例。测试入口为 \path{tests/reference/test_core_monograph.py}。

2026 年 10 月 15 日的 CPU 核验中，该入口的 46 个参数化测试全部通过；仓库全套测试为 316 项通过。积分目标采用 96 点 Gauss--Legendre 求积与解析梯度对照，主要恒等式容差为 $10^{-11}$ 或 $10^{-12}$。U-统计量使用三动作词表、$M=2,3,4$ 的全部有序样本组合；无重复抽样直接比较协方差矩阵，检查半正定差值。局部三次率采用逐渐接近教师的固定方向，不以有限差分噪声替代极限系数。运行环境与命令保存于 \path{validation.json}。

数值测试用于暴露符号、尺度和实现错误；有限实例通过不构成定理证明。理论正文中的结论以相应解析证明和显式条件为依据。测试结果不用于声称共享参数大模型的训练收敛或任务准确率。
